% CLAUDE REPORT::WCD-CONSOLIDATE-0002::2026-11-14T10:50Z
\documentclass[11pt]{article}
\usepackage{amsmath,amssymb}
\usepackage[margin=1in]{geometry}
\title{WCD-CONSOLIDATE-0002: second consolidation of the worst-case derivation routine}
\author{Claude scheduled routine (not an LSPV/TMR role; owns no region of root \texttt{main.tex})}
\date{2026-11-14}
\begin{document}
\maketitle

\section*{Target}
$\overline W_R(N)=\max_{F\in\mathrm{UNSAT},\,|\langle F\rangle|\le N}\ \min_{\pi:F\vdash_R\bot}\mathrm{cost}_R(\pi)$
and the bridge $B_R:\ \mathrm{SAT}\in\mathrm P\Rightarrow\overline W_R(N)\le\mathrm{poly}(N)$.
\textbf{P versus NP remains OPEN.}

\section*{What this run did}
This run claims no new theorem. A survey of remote branches found at least 154 run branches
of this routine (\texttt{ccr-1695bd63-*}, 2026-11-02 to 2026-11-14), none of them merged. The
first consolidation (branch \texttt{vycmx6}, 2026-11-12) was not merged either, so the 31 runs
after it started without it. More than half of those runs re-derived the bridge-collapse
theorem
\[
\mathrm{LB}(R)\ \text{known}\ \Longrightarrow\ \bigl(B_R\iff \mathrm{SAT}\notin\mathrm P\bigr)
\]
or the black-box/oracle collapse, each under a new identifier. This run did three things:
\begin{enumerate}
\item copied the \texttt{vycmx6} consolidated log onto this branch;
\item merged the 31 later runs into stable-state items 7--13 of
  \texttt{CONSOLIDATED\_LIVING\_LOG.md}, with an attribution to a branch for each item, and
  extended the branch index;
\item added a single next target that is not a bridge-equivalence statement.
\end{enumerate}

\section*{CORPUS REVIEWED}
\begin{itemize}
\item \texttt{AGENTS.md}.
\item The current TMR-1 handoff. The active five-agent team in \texttt{main} is TMR, under the 2026-11-13 12:49 directive.
\item The logs and report subjects of all 31 post-consolidation branches.
\item The full consolidated log from \texttt{vycmx6}.
\item The codex reports of 2026-11-11 that define $W_R(N)$.
\end{itemize}
Root \texttt{main.tex} is not edited: this routine has no ownership region in it.

\section*{DUPLICATION / CONTRADICTION AUDIT}
No contradictions were found between runs. The duplication is severe, and its cause is
procedural (unmerged branches), not mathematical. All stable items keep the status their
original report gave them. None is upgraded here.

\section*{Literature this run}
The entries below come from search summaries only; the primary texts were not read this run.
\begin{itemize}
\item ECCC TR26-018 (Feb.\ 2026): Res($\oplus$) refutations of depth $o(n^2/\log^4 n)$ must have
  superpolynomial size.
\item arXiv:2507.23008: the same conclusion at depth $N^{2-\varepsilon}$.
\item DAG-like Res($\oplus$) with no depth cap remains OPEN.
\item arXiv:2605.04544: an IPS-fragment lower bound for CNFs.
\item ECCC TR26-133: lower bounds for PCR over $\mathbb F_2$.
\item CCC 2026: AC$^0[p]$-Frege lower bounds for CNFs remain OPEN for every $p>1$.
\end{itemize}
Classification: each is a \emph{useful analogy needing a transfer theorem}. None reaches a
system that p-simulates the multisign merge calculus extended with unrestricted parity
reasoning, so the load-bearing gap (stable item 6) is unchanged.

\section*{STALE OR OVERSTATED CLAIMS FOUND}
None among the merged items. Item 13 (the survival matrix) is a CANDIDATE and should be
checked against primary sources before load-bearing use.

\section*{Next target}
Read TR26-018 and arXiv:2605.04544 in full. Decide whether either gives a lower bound with
no depth cap for a system that contains the multisign merge calculus plus parity. If not,
name the hypothesis that blocks the transfer.

\section*{Process recommendation}
Merge this branch, or give the routine one stable write branch. Otherwise every run starts
from a fresh clone of \texttt{main} and cannot see earlier results.

% END_OF_VERIFIED_REPORT::CLAUDE::WCD-CONSOLIDATE-0002
\end{document}
